%%%%%%%%%%%%%%%%%%%%%%%%%%%%%%%%%%%%%%%%%%%%%%%%%%%%%%%%%%%%%%%%%%%%%%%%
\section{Near-splits}\label{sec:nearsplit}
%%%%%%%%%%%%%%%%%%%%%%%%%%%%%%%%%%%%%%%%%%%%%%%%%%%%%%%%%%%%%%%%%%%%%%%%

\subsection{Thin necks in finite meanders}

A meander of order $2m$ has a \emph{thin neck at the middle} if the first $2m$ of its $4m$
crossings along the river occupy $2m$ consecutive positions along the road, read cyclically;
equivalently, exactly two road arches pass over or under the gap between crossings $2m$ and
$2m+1$. A vertical line through that gap then meets the road twice, and on the sphere it is a short
closed curve cutting the surface into equal halves. Call this the \emph{zero-defect event} $F_0$:
no crossing is out of place. Write $\Prob_{2m}$ for the uniform probability on meanders of order
$2m$. Call a meander of order $2m$ a \emph{junction} if it is two meanders of order $m$ glued end
to end (Section~\ref{sec:intro}); only two arches pass over its middle gap, the short arch
joining crossings $2m$ and $2m+1$ and the long arch over it, so it is a special case of $F_0$. Since gluing
is injective, $q_m=M_m^2/M_{2m}$ is the probability that a uniform meander of order $2m$ is a
junction. Let $X_m$
be the number of meanders of order $m$ with one marked \emph{exterior} arch above the river, an
arch not nested under any other arch. Then
\begin{equation}\label{eq:F0}
\Prob_{2m}(F_0)=\frac{2X_m^2}{M_{2m}},\qquad M_m\le X_m\le\tfrac12M_{m+1},\qquad
2q_m\le\Prob_{2m}(F_0)\le\tfrac12C_r^2\,q_m ,
\end{equation}
with $C_r=\sup_nM_{n+1}/M_n<\infty$. (By parity the two arches are both above or both below the
river; given the side, meanders with $F_0$ correspond to pairs of meanders with a marked exterior
arch.) Under~\eqref{eq:conj},
\[
q_m\approx\frac{(A\mu^mm^{-\alpha})^2}{A\mu^{2m}(2m)^{-\alpha}}=A\,2^\alpha m^{-\alpha},
\]
so $\alpha$ is the exponent of a thin neck cutting a meander into equal halves~\cite{PaperA}.

\subsection{The continuum event}

Let $\Pi$ be a permuton, and read its second coordinate, the road, on the circle
$\mathbb T=\mathbb R/\mathbb Z$, since the road is a loop. Let $I=[0,1/2]$ be the first half of the river and $J_s=[s,s+1/2]$
(mod $1$) a cyclic half of the road. The \emph{off-block mass}
\[
d_\Pi(s)=\Pi\bigl(I\times(\mathbb T\setminus J_s)\bigr)+\Pi\bigl(([0,1]\setminus I)\times J_s\bigr)
\]
counts crossings in the first river half but not in $J_s$, or the other way around.
Intuitively, it measures how far the crossings are from splitting into two blocks at the cut
points $1/2$ on the river and $s$, $s+\frac12$ on the road. For the meandric permuton,
\[
d_\Pi(s)=\nu\bigl(\eta_1([0,1/2])\,\triangle\,\eta_2(J_s)\bigr),
\]
the area of the symmetric difference of the regions covered by the river in its first half of time
and by the road during $J_s$. A \emph{near-split with tolerance $\delta$} means
$\inf_sd_\Pi(s)\le\delta$. For finite meanders it is convenient to ignore thin strips around the cut
points. For
$0<\delta<1/8$, let $J_s^{+\delta}$ be the closed $\delta$-neighborhood of $J_s$ and $J_s^{-\delta}$
the set of points of $J_s$ at distance at least $\delta$ from its complement, and put
\[
d^\delta_\Pi(s)=\Pi\bigl([0,\tfrac12-\delta]\times(\mathbb T\setminus J_s^{+\delta})\bigr)
+\Pi\bigl([\tfrac12+\delta,1]\times J_s^{-\delta}\bigr),\qquad
E_\delta=\Bigl\{\inf_sd^\delta_\Pi(s)\le\delta\Bigr\}.
\]
For a permuton the two versions differ by at most $6\delta$, since by the uniform marginals a strip
of width $\delta$ has mass at most $\delta$. Hence
\begin{equation}\label{eq:sandwich}
\Bigl\{\inf_sd_\Pi(s)\le\delta\Bigr\}\subset E_\delta\subset\Bigl\{\inf_sd_\Pi(s)\le7\delta\Bigr\},
\end{equation}
and bounds on the two events are interchangeable up to a constant change of $\delta$.

\subsection{The exact finite version}

For a meander of order $n$, the \emph{centered atomic plot} $\Pi_n$ has mass $1/(2n)$ at each point
$\bigl((i-\frac12)/2n,(\omega(i)-\frac12)/2n\bigr)$; it has the same limits as the permuton.

\begin{proposition}\label{prop:micro}
For a meander of order $n$ and $0<\delta<1/(4n)$, the event $E_\delta$ for $\Pi_n$ holds exactly
when the first $n$ crossings along the river occupy $n$ consecutive positions along the road, read
cyclically. In particular, for order $2m$ and $\delta<1/(8m)$ it is the zero-defect event $F_0$,
and its probability is given by~\eqref{eq:F0}.
\end{proposition}

\begin{proof}
Each atom has mass $1/(2n)>\delta$, so $d^\delta_{\Pi_n}(s)\le\delta$ forces $d^\delta_{\Pi_n}(s)=0$. No
atom lies within $\delta$ of the river midpoint. The endpoints of $J_s$ are exactly $1/2$ apart, so
they sit alike relative to the grid of road positions, which has spacing $1/(2n)$. If neither is within $\delta$ of an
atom, $J_s^{+\delta}$ and $J_s^{-\delta}$ contain the same $n$ consecutive positions, and zero
off-block mass says that these are the first-half crossings. Otherwise $J_s^{-\delta}$ contains $n-1$
consecutive positions, all of first-half crossings, and $J_s^{+\delta}$ adds one candidate at each end;
exactly one of the two candidates is a first-half crossing, which again gives $n$ consecutive
positions. Conversely, if the first-half crossings occupy $n$ consecutive positions, take $s$ at the
midpoint between grid positions just before them; both endpoints of $J_s$ are then at distance
$1/(4n)>\delta$ from all atoms, and the off-block mass is zero.
\end{proof}

Heuristically, at resolution $\delta$ a meander looks like one of order about $1/\delta$, so by the
estimate $q_m\approx A2^\alpha m^{-\alpha}$ one expects $\Prob(E_\delta)\approx\delta^{\astar}$ in the
limit.

\subsection{The bounds}

\begin{theorem}[near-splits in the conjectured limit]\label{thm:nearsplit}
Let $\Pi$ be the meandric permuton of Section~\ref{sec:limit}, with
$\gamma^2=(17-\sqrt{145})/3$. There are constants $0<c_1\le C_1$ such that for all small $\delta>0$
\[
c_1\,\delta^{\astar}\;\le\;\Prob(\Pi\in E_\delta)\;\le\;C_1\,\delta^{\xstar+5/7},
\qquad\astar=\frac{29+\sqrt{145}}{12}\approx3.4201,\quad\xstar+\frac57\approx3.1344 .
\]
\end{theorem}

This concerns the limit object only; Section~\ref{sec:bridge} discusses finite meanders. The lower
bound has the predicted exponent, and holds for every
$\gamma\in(0,\sqrt2)$, that is, for every $x=4/\gamma^2>2$, with exponent $1+x$. The exponent of the upper bound is the best we
have reached. If the limit $\lim_{\delta\to0}\log\Prob(E_\delta)/\log\delta$ exists, it lies in
$[3.1344,3.4201]$. Two side results give context.
\begin{itemize}
\item \emph{Exact splits do not occur.} Almost surely no region covered by the river in a proper
cyclic time interval equals, up to area zero, one covered by the road; so $\Prob(E_0)=0$. (The
boundary of the region an SLE$_8$ has covered by a given time consists of SLE$_2$-type curves, and
the boundaries for two independent curves share no arc.)
\item \emph{One configuration has the sharp exponent.} The field event that two fixed disks around
the marked points $0$ and $1$ of the sphere (see below) each carry area $1/2+O(\delta)$, and the rest
$O(\delta)$, has probability $\delta^{\astar+o(1)}$.
\end{itemize}

\subsection{The lower bound}

The lower bound is proved by building the event, with four ingredients.
\emph{Two blobs suffice.} Suppose the sphere splits, up to area zero, into three regions, two
\emph{blobs} $O_1$, $O_2$ and a remainder $O_3$, with $|\nu(O_1)-1/2|\le\delta_1$ and $\nu(O_3)\le\delta_2$, and suppose each curve, after
discarding the times it spends in $O_3$, visits all of one of $O_1,O_2$ before any of the other.
Then $\inf_sd_\Pi(s)\le2\delta_1+7\delta_2$, and so $E_{2\delta_1+7\delta_2}$ holds. This is deterministic:
if both curves finish one blob before starting the other, the first half of river time and a
suitable half of road time cover roughly the same blob.

\emph{Fixed clumps are visited in blocks.} Fix two points $z_1,z_2$ and small disks of radius $r$
around them. Say a curve visits the two disks \emph{in blocks} if it fills one disk completely
before entering the other. Each curve visits $z_1$ and $z_2$ once, at different times, so by
continuity, for small $r$, it visits the two disks in blocks. Both curves do so with some probability $P_{\mathrm{blk}}(r)>0$; as the curves
are independent of the field and the area parametrization changes speeds but not orders, this holds
conditionally on any field event.

\emph{A convenient normalization of the field.} The unit-area sphere comes from a field with additive
constant fixed on the unit circle and three $\gamma$-singularities, at the marked points $0$, $1$
and $\infty$: the terms $-\gamma\log|z|$, $-\gamma\log|z-1|$ and, for $|z|>1$, $-(2Q-\gamma)\log|z|$
with $Q=2/\gamma+\gamma/2$~\cite[Definition~2.2]{BGS}. Its law is weighted by the power
$(2Q-3\gamma)/\gamma=x-2$ of its total area, where $x=4/\gamma^2$, and then normalized; the range
$\gamma\in(0,\sqrt2)$ of the lower bound is exactly $x>2$, where this power is positive. A Gaussian
change of variables rewrites this law exactly with the constant fixed on a large circle, inside
which the term from $\infty$ is constant. So near $0$ and $1$ the field is a Gaussian free field
plus the two singularities and a constant, and its restrictions to small disks around $0$ and $1$
are conditionally independent given the outside.

\emph{Pinned clumps, balance and the sphere weight.} Call a disk \emph{pinned} if it is centered at
a marked point, where the field has the singularity $-\gamma\log|z-z_0|$. The area of a pinned disk
has a heavier tail than that of an ordinary disk: by Wong's analysis of such chaos
masses~\cite{Wong}, it lies in $[t,2t]$ with probability asymptotic to a constant times
$t^{-(x-1)}$, against $t^{-x}$ for an ordinary disk. (The marked points are typical points of the
area, where the field is already high.) Given the outside, the disks at $0$ and $1$ are independent,
so both reach area of order $t$, with bounded outside area, at probability of order $t^{-2(x-1)}$. A
one-dimensional Gaussian coordinate of the field in the first disk moves its area continuously; the
two areas agree within $1$ with conditional probability of order $1/t$. Then after normalization
each disk holds area $1/2+O(1/t)$ and the rest $O(1/t)$. Finally the weight $(\text{total
area})^{x-2}$ multiplies the probability by $t^{x-2}$. Altogether
\[
t^{-2(x-1)}\cdot t^{-1}\cdot t^{\,x-2}=t^{-(1+x)} ,
\]
and with $\delta$ of order $1/t$ the first two ingredients give $\Prob(E_\delta)\ge c_1\delta^{1+x}$,
which is $\delta^{\astar}$ at $x=\xstar$. The disks have a fixed radius with $P_{\mathrm{blk}}(r)>0$; no estimate
for shrinking disks is needed.

\subsection{The upper bound}

The upper bound splits the event $\inf_sd_\Pi(s)\le\delta$, equivalent by~\eqref{eq:sandwich},
by where the area sits.
\begin{itemize}
\item \emph{Concentrated area.} If one small ball holds almost all the area, or the boundary of the
region the river covers in its first half of time is very small, the probability is at most a fixed fraction of that of the whole
event. This uses re-rooting: moving the marked points to points sampled from the area. These cases
are absorbed into the others.
\item \emph{Diffuse area with a moderate total.} If the total area $t$ before normalization is at
most $\delta^{-\vartheta}$ for a fixed $\vartheta<1$, and the boundary is macroscopic, near-agreement forces the boundaries of
the two regions to come close at arbitrarily many separated places, which for independent
SLE$_2$-type curves has probability below any power of $\delta$, by multipoint SLE estimates.
\item \emph{A large total.} If $t>\delta^{-\vartheta}$, two separated cells of a fixed mesh carry fixed
fractions of it. The worst case is two cells at the marked points, costing
$t^{-2(x-1)}$, refunded by $t^{x-2}$ from the sphere weight: net $t^{-x}$, at most $\delta^{\vartheta x}$,
with $\vartheta$ as close to one as desired.
\end{itemize}
So $\Prob(E_\delta)\le C\delta^{x-o(1)}$, one power short: the missing power is the balance factor of
the lower bound. The five-sevenths corollary cited in the appendix recovers part of this power through estimates on
retained rank integrals and a sum over configurations. This gives $x+5/7$. We do not interpret that integral estimate
as a conditional probability of visit-order agreement. The prediction $1+x=\astar$ corresponds to recovering the full factor.
